% Part: normal-modal-logic
% Chapter: filtrations
% Section: S5-fmp

\documentclass[../../../include/open-logic-section]{subfiles}

\begin{document}

\olfileid[id]{nml}{fil}{fmp}

\olsection{\Log{K} dan \Log{S5} Mempunyai Sifat Model Berhingga}

\begin{defn}
  Suatu sistem logika modal $\Sigma$ dikatakan mempunyai \emph{sifat model
    berhingga} jika, setiap kali !!a{formula}~$!A$ bernilai benar pada suatu
  dunia dalam suatu model dari $\Sigma$, $!A$ juga bernilai benar pada suatu
  dunia dalam suatu model \emph{berhingga} dari~$\Sigma$.
\end{defn}

\begin{prop}\ollabel{prop:K-fmp}
  \Log{K} mempunyai sifat model berhingga.
\end{prop}

\begin{proof}
  \Log{K} merupakan himpunan !!{formula}s yang valid, yakni setiap model
  merupakan model dari~\Log{K}. Menurut \olref[fil]{thm:filtrations}, jika
  $\mSat{M}{!A}[w]$, maka $\mSat{M^*}{!A}[{[w]}]$ bagi setiap filtrasi
  dari~$\mModel{M}$ melalui himpunan $\Gamma$ semua sub-!!{formula}s
  dari~$!A$. Setiap !!{formula} hanya mempunyai berhingga banyak
  sub-!!{formula}s, sehingga $\Gamma$ berhingga. Menurut
  \olref[fin]{prop:filt-are-finite}, $\card{W^*} \le 2^n$, dengan $n$
  banyaknya !!{formula}s dalam~$\Gamma$. Karena \Log{K} tidak memberlakukan
  pembatasan apa pun pada model, $\mModel{M^*}$ merupakan model dari~\Log{K}.
\end{proof}

Untuk menunjukkan melalui filtrasi bahwa suatu logika~\Log{L} mempunyai sifat
model berhingga, sangat penting bahwa filtrasi dari suatu model dari~\Log{L}
itu sendiri merupakan model dari~\Log{L}. Hal ini sering memerlukan cukup
banyak upaya, dan tidak setiap filtrasi menghasilkan model dari~\Log{L}. Namun, bagi model
universal, hal itu tetap berlaku.

\begin{prop}\ollabel{prop:univ-fin}
  Misalkan $\mClass{U}$ merupakan kelas model universal (lihat
  \olref[frd][es5]{prop:S5=univ}) dan $\mClass{U}_\mathrm{Fin}$ merupakan
  kelas semua model universal berhingga. Maka sebarang !!{formula}~$!A$ valid
  dalam $\mClass{U}$ jika dan hanya jika formula itu valid dalam
  $\mClass{U}_\mathrm{Fin}$.
\end{prop}

\begin{proof}
  Model universal berhingga merupakan model universal, sehingga arah dari
  kiri ke kanan langsung berlaku. Untuk arah dari kanan ke kiri, andaikan
  bahwa~$!A$ bernilai salah pada suatu dunia $w$ dalam suatu model universal
  $\mModel{M}$. Misalkan $\Gamma$ memuat $!A$ beserta semua subformulanya;
  jelas bahwa $\Gamma$ berhingga. Ambil suatu filtrasi $\mModel{M^*}$
  dari~$\mModel{M}$; maka $\mModel{M^*}$ berhingga menurut
  \olref[fin]{prop:filt-are-finite}, dan menurut
  \olref[fil]{thm:filtrations}, $!A$ bernilai salah pada $[w]$
  dalam~$\mModel{M^*}$. Tinggal diamati bahwa $\mModel{M^*}$ juga universal:
  jika diberikan $u$ dan $v$, berdasarkan hipotesis berlaku $Ruv$, dan
  berdasarkan \olref[fil]{defn:filtration}\olref[fil]{defn:filtration-R},
  berlaku pula $R^*[u][v]$.
\end{proof}

\begin{cor}\ollabel{cor:S5fmp}
  \Log{S5} mempunyai sifat model berhingga.
\end{cor}

\begin{proof}
  Menurut \olref[frd][es5]{prop:S5=univ}, jika $!A$ bernilai benar pada suatu
  dunia dalam suatu model refleksif dan Euklidean, maka formula itu bernilai
  benar pada suatu dunia dalam suatu model universal. Menurut
  \olref{prop:univ-fin}, formula itu bernilai benar pada suatu dunia dalam
  suatu model universal berhingga (yakni filtrasi model tersebut melalui
  himpunan sub-!!{formula}s dari~$!A$). Setiap model universal juga refleksif
  dan Euklidean; jadi, $!A$ bernilai benar pada suatu dunia dalam suatu model
  refleksif dan Euklidean berhingga.
\end{proof}

\begin{prob}
  Tunjukkan bahwa setiap filtrasi dari suatu model serial atau refleksif juga
  serial atau refleksif (masing-masing).
\end{prob}

\begin{prob}
  Temukan suatu filtrasi tak simetris (tak transitif, tak Euklidean) dari suatu
  model simetris (transitif, Euklidean).
\end{prob}

\end{document}
